\section{An exact renewal envelope from the first missing weight}
Let $q\ge2$, $S_q=z+\cdots+z^q$, and define
\[
 H_q(z)=\frac1{1-S_q(z)},\qquad
 G_q(z)=\frac1{(1-z)(1-S_q(z))},\qquad
 b_q(z)=(1-z^{q-1})G_q(z)=\frac{[q-1]}{1-S_q(z)}.
\]
Thus $H_q$ counts all words with parts at most $q$, and $G_q$
is its forward cumulative sequence. All coefficients of $G_q$ are
positive integers and nondecreasing.

\begin{lemma}[Log-concavity of the renewal cumulative]
The sequence $G=G_q$ is log-concave at every rank.
\end{lemma}
\begin{proof}
Its initial values are $G_n=2^n$ for $0\le n\le q$, followed by
$G_{q+1}=2^{q+1}-1$. The generating function gives
\[
                 G_{n+1}=2G_n-G_{n-q}\qquad(n\ge q).
\]
The first $q-1$ log-concavity minors vanish and the rank-$q$ minor
is $2^{q-1}$. At every $n\ge q+1$ the minor equals
\[
 G_n^2-G_{n-1}G_{n+1}
       =G_{n-1}G_{n-q}-G_nG_{n-q-1}.
\]
By induction the earlier adjacent ratios are nonincreasing, so
$G_{n-q}/G_{n-q-1}\ge G_n/G_{n-1}$. This pays the displayed
minor and proves the unbounded induction. No log-concavity of
the actual arm's forward cumulative sequence is assumed.
\end{proof}

\begin{corollary}[Own-cumulative C for the untruncated alphabet reference]
For every $q\ge2$, the infinite word enumerator $H_q$ itself
satisfies property C at every shift $r\ge3$.
This is a statement about the infinite alphabet-restricted
reference, not every finite composition host.
\end{corollary}
\begin{proof}
Its shifted positive derivative is
$d_j=[z^j][q-1]H_q=G_q(j)-G_q(j-q+1)$.
Put $k=q-1\ge1$ and abbreviate $G=G_q$, with zero extension.
For $j\ge1$, its shift-three C minor has the exact decomposition
\[
\begin{split}
 G_jd_j-G_{j-1}d_{j+1}
 &=\{G_j^2-G_{j-1}G_{j+1}\}\\
 &\quad+\{G_{j-1}G_{j+1-k}-G_jG_{j-k}\}\ge0.
\end{split}
\]
The first bracket is the already proved cumulative-LC minor.
The second is nonnegative by the same decreasing adjacent
ratios, with earlier zero indices handled directly.
When $k=1$ it vanishes identically.
The initial rank $j=0$ is immediate. Finally cumulative
log-concavity makes the earlier denominator ratio for every
$r>3$ at least the shift-three ratio, proving all shifts.
The generic ratio-chain argument is already kernel checked
in \texttt{FirstMissingRenewalEnvelope.lean}; no new separate
end-to-end renewal/composition formalization is claimed here.
\end{proof}

\begin{theorem}[First-missing relative envelope]
Let $X$ be an actual nonrectangular increasing arm, let
$\kappa\ge3$ be its first missing composition weight, and put
$q=\kappa-1$. Then, with
$d_n=((F_X)_{n+2}-(F_X)_{n+1})^+$,
\[
                         \frac{d_n}{G_q(n)}
                    \quad\text{is nonincreasing in }n.       \tag{1}
\]
The same conclusion holds on any larger formal arm class on which
the already proved second-positive-difference source bound and
the ordinary positive-difference unimodality hold.
In particular, $d_n\le G_q(n)$ and its positive adjacent ratios
are bounded by the explicit decreasing ratios $G_q(n)/G_q(n-1)$.
These tend to the unique positive root $\rho_q$ of
$\rho_q^{-1}+\cdots+\rho_q^{-q}=1$.
For $\kappa=3$, $G_2(n)=F_{n+3}-1$, giving an envelope that does
not depend on any capacity count or upper bound.
\end{theorem}
\begin{proof}
Set $w_n=(d_n-d_{n-1})^+$, with $d_{-1}=0$, and
$a_n=\sum_{i\le n}w_i$. By the proved shifted second-difference
source theorem,
\[
            v=(1-z^{q-1})-(1-S_q)w\quad\text{has }v_i\ge0.
\]
Dividing this exact formal-power-series identity by
$(1-z)(1-S_q)$ gives
\[
                      a=b_q-G_qv.                            \tag{2}
\]
Log-concavity of $G_q$ means that each ratio
$G_q(n-i)/G_q(n)$, with negative indices set to zero, is
nondecreasing in $n$. Moreover
$b_q(n)/G_q(n)=1-G_q(n-q+1)/G_q(n)$ is nonincreasing.
Equation (2), with every $v_i\ge0$, therefore proves that
$a_n/G_q(n)$ is nonincreasing. Every coefficient comparison
involves only a finite sum, so no convergence issue arises.

The already proved unimodality of the positive derivative $d$
means that $a=d$ through its first maximum and $a$ stays constant
afterwards. Up to that maximum use the preceding relative bound;
afterwards $d$ is nonincreasing while $G_q$ is nondecreasing.
This proves (1), including its joining edge. The root limit follows
from the elementary finite-order positive renewal recurrence.
\end{proof}

\paragraph{Exact scope and a new obstruction.}
This comparison does not prove the own-cumulative property C.
Indeed take $q=2$ and the nonnegative single renewal loss $v=z^3$.
The nondecreasing integer sequence
\[
       d_n=b_2(n)-G_2(n-3)
\]
has nonnegative renewal-source loss exactly $z^3$. Reconstruct
$F_0=F_1=1$, $F_{n+2}=F_{n+1}+d_n$, and $C_n=\sum_{i\le n}F_i$.
At $n=8$,
\[
                    d_8C_6-d_7C_7=19>0,
\]
so $d_n/C_{n-1}$ increases. This is an abstract source-admissible
sequence, not a claimed actual arm. It proves that the new
second-difference source and renewal envelope alone cannot close
property C. Actual word coupling must still be retained.
General original unamplified coefficient unimodality remains OPEN.
